\section{Related Work}
\label{sec:rw}

The landscape of task scheduling in data centers keeps evolving. It has moved from traditional prioritization towards much more intricate online learning frameworks. Traditional scheduling research was more focused on singular objectives like makespan, using algorithms like HEFT and PEFT \citep{HosseiniShirvani2021A}. These traditional algorithms are heavily reliant on the fact that the physical hardware is statically available; it does not factor in downtime into its scheduling process. This makes it suffer serious performance degradation when it experiences hardware faults. To help with these limitations, newer algorithms like Hybrid Swarm Optimizers were developed \citep{Alkaam2025Hybrid}.

The latest advancements in the space seem to have moved towards utilizing Deep Reinforcement Learning and Multi-Armed Bandit Algorithms \citep{Khan2025Machine,DeCarvalho2025Toward}. Using the Deep Reinforcement Learning methodology provides stellar results in stationary environments, but have very high training overhead \citep{Choppara2025Efficient,Sravan2026DRL}. Multi-Armed Bandit gives the option of using hyper-heuristics for lightweight online learning \citep{Singh2024Leveraging}

There are still three major gaps in the current literature:
\begin{enumerate}
    \item \textbf{Historical Bias in Non-Stationary Environments:} Standard MAB algorithms such as UCB and standard Thompson Sampling keep a very long performance history window, which makes it difficult for them to quickly react to unexpected capacity drops, or other hardware failures.
    \item \textbf{Lack of Lightweight Contextual Workload Fingerprinting:} Existing contextual or content-aware schedulers often rely on heavy offline classifiers or neural feature extractors \citep{Adil2022scp}. 
    \item \textbf{Unbalanced Multi-Objective Reward Formulations:} While sustainability and carbon intensity have emerged as first-class scheduling parameters alongside cost and latency \citep{Arora2023Towards}, existing multi-objective reward functions often suffer from scale distortion across heterogeneous execution batches, leading to single-metric dominance during bandit update steps.
\end{enumerate}

The remainder of this section critically reviews prior contributions across these foundational domains.

\subsection{Foundations of Non-Stationary Workload Allocation}
In heterogeneous cloud environments where tasks and requests don't all arrive in a sequential or predictable order, allocating resources becomes an NP-hard problem.

\citet{Mathew2014Study} categorised different scheduling algorithms around performance measures such as makespan, load balancing and throughput. \citet{Saurav2021A} takes this further by looking into energy-aware task scheduling in large, heterogeneous systems, and bargaining and exchange between execution speed and power consumption. All literature referenced in this section come to the same conclusion that traditional scheduling algorithms always underperform in heterogeneous environments.

\subsection{Deterministic Task Prioritization and Heuristic Approximations}
List-based algorithms like HEFT, PEFT and HSIP work in two stages: firstly, they rank the tasks, then they assign them to processors. This flow helps to make DAG scheduling less complex, and it provides predictable scheduling behaviour.

\citet{HosseiniShirvani2021A} uses this approach as a foundation to propose HH-LiSch: a hybrid list scheduler that combine two ranking strategies RankI and RankO. It was tested on benchmarks like Montage and Cybershake, and it saw a reduction in average makespan by 16.71\%, compared to HEFT and PEFT.

The major limitation is still the fact that classical list schedulers use fixed rules and limited workload information, and they do not learn from changing task characteristics.

\subsection{Content and Workload Feature Engineering for Scheduling}
Instead of treating each task as a singular unit of computation, using feature-informed schedulers use properties of the workload itself to examine task characteristics and determine resource allocation.

\citet{Adil2022scp} proposed CA-MLBS, a solution that combines both Support Vector Machine classification and Particle Swarm Optimization. The algorithm works by classifying incoming tasks by content type and maps them onto the most ideal machine available. Experiments showed that this method reduced makespan by up to 29\% and increased throughput by 44\% against feature-agnostic algorithms.

A major deterrent with this approach is that the classifier is trained offline, which means that once it is deployed, it does not adapt to changes in hardware or workload conditions.

\subsection{Complex Search Space Exploration in Hybrid Meta-Heuristics}

\citet{Kumar2021Implementation} applied Ant Colony Optimization (ACO) in CloudSim and reported lower average task execution times and cloudlet failure rates than Round Robin and Minimum Completion Time (MCT) policies.

Recent studies have gone further by combining several meta-heuristics, like \citet{Alkaam2025Hybrid}, who proposed HGHHC: which combines Henry Gas Solubility Optimization, Harris Hawks Optimization, and Comprehensive Opposition-Based Learning (COBL). Across HPC2N, NASA, and synthetic workload traces, HGHHC reduced makespan by 34.30\%, 72.95\%, and 28.67\% respectively, relative to the reported baseline swarm methods. \citet{Bal2022A} also proposed RATS-HM, which combines an improved cat swarm optimization scheduler (ICS-TS), a group optimization deep neural network (GO-DNN) and NSUPREME encryption for resource allocation and security.

The downside of this is that it is very computationally expensive, as these methods most likely require multiple population-based searched to properly assign tasks to resources. Having to repeatedly run this process for constantly changing workloads is expensive and may as well be too slow.

\subsection{Adaptive Online Decision-Making and Regret Minimization}

Multi-Armed Bandit models continuously learns from previous actions, rather than searching from scratch. The trade-off is between exploration and exploitation: where the algorithm tries lesser known actions, or it continues with actions that have previously performed well.

\citet{Singh2024Leveraging} proposed a MAB model for task scheduling that derives logarithmic regret bounds, as well as proposes latency and utilization improvements. \citet{Tran-Dang2024Parallel} merged Thompson Sampling with matching theory for a decentralised taks offloading through MAB-MATCH. \citet{Chigu2024A} developed an Upper Confidence Bound scheduler for serverless FaaS systems. Multi-Armed Bandits can also be used for federated learning under mobility constraints, as well as satellite beamforming \citep{Jun2025Mobility,Liang2026Joint}.

This approach stands out from the rest because of how highly adaptable it is. MAB-based algorithms can easily and confidently update their decisions as new information and performance feedback arrive. This is why they are a standout candidate for dynamic workloads.

A detailed taxonomy and feature comparison summarizing the literature alongside the characteristics of our proposed contextual MAB hyper-heuristic framework is presented in Table~\ref{tab:related_work_summary}.

\begin{table}[H]
\centering
\caption{Comprehensive Comparison of Literature and Features of the Proposed Work}
\label{tab:related_work_summary}
\small
\resizebox{\textwidth}{!}{%
\begin{tabular}{p{3.2cm} p{2.8cm} p{2.5cm} p{2.8cm} p{2.5cm} p{2.2cm} p{2.2cm}}
\toprule
\textbf{Study / Work} & \textbf{Methodology / Core Strategy} & \textbf{Optimization Objectives} & \textbf{Target Environment} & \textbf{Dynamic Adaptivity} & \textbf{Context / Feature Awareness} & \textbf{Evaluation Engine} \\
\midrule

\multicolumn{7}{l}{\textbf{\textit{List-Based \& Classical Heuristics}}} \\
\citet{HosseiniShirvani2021A} & HH-LiSch (Task duplication \& slot insertion) & Makespan & Cloud VMs & Static / Low & $\times$ & Simulation \\
\citet{Adil2022scp} & CA-MLBS (SVM + PSO content classification) & Makespan, Throughput & Cloud & Static & $\checkmark$ (Content Type) & CloudSim \\

\midrule
\multicolumn{7}{l}{\textbf{\textit{Meta-Heuristic \& Hybrid Approaches}}} \\
\citet{Kumar2021Implementation} & Ant Colony Optimization (ACO) & Execution Time, Failure Rate & Cloud & Moderate & $\times$ & CloudSim \\
\citet{Alkaam2025Hybrid} & HGHHC (Henry Gas + Harris Hawks + COBL) & Makespan & Heterogeneous Cloud & Moderate & $\times$ & HPC2N / NASA Traces \\
\citet{Bal2022A} & RATS-HM (Cat Swarm + GO-DNN + Encryption) & Utilization, Power, Security, Latency & Cloud / Edge & Moderate & $\times$ & CloudSim \\

\midrule
\multicolumn{7}{l}{\textbf{\textit{Machine Learning \& Neural Network Methods}}} \\
\citet{Daid2021An} & Probabilistic Neural Networks + Linear Regression & CPU Power, Allocation & Cloud Datacenters & High & $\checkmark$ (CPU Load) & MATLAB (10 VMs) \\
\citet{Sravan2026DRL} & Multi-Objective DRL Framework & Latency, Energy, SLA & Cloud / Fog & High & $\checkmark$ (SLA Bounds) & Simulation \\

\midrule
\multicolumn{7}{l}{\textbf{\textit{Multi-Armed Bandit (MAB) \& Hyper-Heuristic Schedulers}}} \\
\citet{Singh2024Leveraging} & Game-Theoretic MAB & Utilization, Latency, Fairness & Dynamic Cloud & High & $\times$ & Cloud Simulator \\
\citet{DeCarvalho2025Toward} & HHCSP (UCB-MAB + FRRMAB Credit Assignment) & Makespan, Cost, Carbon & Cloud (AWS / Teads) & High & $\times$ & Real AWS Data \\
\citet{Chigu2024A} & UCB-based MAB Scheduler & Execution Latency & Edge-Cloud FaaS & High & $\times$ & AWS Lambda / OpenWhisk \\



\midrule
\rowcolor{gray!15}
\textbf{Proposed Work} & \textbf{Contextual MAB Hyper-Heuristic (UCB / Discounted-UCB / Thompson) + IPC Bridge} & \textbf{Makespan, Utilization, Non-stationary Regret} & \textbf{Heterogeneous Edge-Cloud} & \textbf{High (Online Adaptive)} & \textbf{$\checkmark$ (Task Graph \& Workload Features)} & \textbf{CloudSim Plus (Java) + Python IPC Bridge} \\
\bottomrule
\end{tabular}%
}
\end{table}

